\section*{第 \ref{ch:g12:space} 章 --- 空间中的向量、直线与平面}

\begin{solution}{exo:g12:space:1}
$\vect{AB}(2, 1, -1)$ 与 $\vect{AC}(1, -1, 1)$ 不成比例
（$\frac21 \neq \frac{1}{-1}$），故三点不共线。直线 $(AB)$ 的方程为：
\[
(x, y, z) = (1 + 2t,\; t,\; 2 - t), \qquad t \in \R .
\]
\end{solution}

\begin{solution}{exo:g12:space:2}
$2(x - 1) - (y - 1) + 3(z - 0) = 0$，\emph{即}
\[
2x - y + 3z - 1 = 0 .
\]
对 $B(0,2,1)$，$0 - 2 + 3 - 1 = 0$，成立，故 $B$ 在该平面上。
\end{solution}

\begin{solution}{exo:g12:space:3}
$\vect{AB}(1,1,0)$，$\vect{AC}(1,0,1)$：
\[
\cos\widehat{A}
= \frac{\vect{AB}\cdot\vect{AC}}{\norm{\vect{AB}}\,\norm{\vect{AC}}}
= \frac{1}{\sqrt2\,\sqrt2} = \frac12,
\]
于是 $\widehat A = 60^\circ$。
\end{solution}

\begin{solution}{exo:g12:space:4}
直线 $\mathcal D$ 上的点为 $(1 + t,\, 2 - t,\, 2t)$。将其代入
$\mathcal P$ 的方程，得：
\[
2(1+t) + (2 - t) - 2t + 1 = 5 - t = 0 \iff t = 5 .
\]
唯一交点为 $(6, -3, 10)$。
\end{solution}

\begin{solution}{exo:g12:space:5}
两平面的法向量 $(1, -1, 2)$ 与 $(2, 1, 1)$ 不共线，故两
平面相交于一条直线。两式相加得 $3x + 3z = 5$，于是
$x = \frac53 - z$；再由第一式得
$y = x + 2z - 1 = \frac23 + z$。记 $z = t$：
\[
(x, y, z) = \left(\frac53 - t,\; \frac23 + t,\; t\right), \qquad t \in \R,
\]
所得为方向向量为 $(-1, 1, 1)$ 的直线。（两方程代入后
恒成立。）
\end{solution}

\begin{solution}{exo:g12:space:6}
两直线的方向向量 $\vec u_1(1, 2, -1)$ 与 $\vec u_2(1, 1, 2)$ 不共线，
故两直线不平行。若二者相交，则需
\[
1 + t = 2 + s, \qquad 2t = 1 + s, \qquad 3 - t = 1 + 2s .
\]
由前两式得 $t = 1 + s$ 及 $2(1+s) = 1 + s$，于是 $s = -1$，$t = 0$；
代入第三式得 $3 = -1$，矛盾。故两直线无公共点，且不平行，
两直线为异面直线。
\end{solution}

\begin{solution}{exo:g12:space:7}
\emph{1.} 距离 $\operatorname{dist} = \dfrac{\abs{2\cdot3 - 1 + 2\cdot1 - 3}}
{\sqrt{4 + 1 + 4}} = \dfrac{4}{3}$.

\emph{2.} 设 $H = M_0 + t\,\vec n$，其中 $\vec n(2, -1, 2)$，选取 $t$
使 $H \in \mathcal P$：
\[
2(3 + 2t) - (1 - t) + 2(1 + 2t) - 3 = 4 + 9t = 0 \iff t = -\frac49 .
\]
故 $H\left(\frac{19}{9}, \frac{13}{9}, \frac19\right)$，且
\[
M_0H = \norm{t\,\vec n} = \frac49 \times 3 = \frac43,
\]
与距离公式的结果一致。
\end{solution}

\begin{solution}{exo:g12:space:8}
坐标：$G(1,1,1)$，平面 $(BDE)$ 经过点 $B(1,0,0)$，
$D(0,1,0)$，$E(0,0,1)$。

\emph{1.} 平面 $(BDE)$ 的方程为 $x + y + z = 1$（三点满足该方程，
且三点不共线），因此 $\vec n(1,1,1)$ 为该平面的
法向量。由 $\vect{AG} = (1,1,1) = \vec n$ 知，对角线 $(AG)$ 与
平面 $(BDE)$ 正交。

\emph{2.} 点 $A(0,0,0)$ 到平面的距离：
$\frac{\abs{0 + 0 + 0 - 1}}{\sqrt3} = \frac{1}{\sqrt3} = \frac{\sqrt3}{3}$.
直线 $(AG)$ 为 $(t, t, t)$；当 $3t = 1$ 时与平面相交，
交点为 $\left(\frac13, \frac13, \frac13\right)$，它是
三角形 $BDE$ 的重心。
\end{solution}

\begin{solution}{exo:g12:space:9}
\emph{1.} $\vect{AB}(1, 0, -1)$，$\vect{AC}(-1, 1, 0)$，
$\vect{AD}(1, 1, 1)$。若 $\vect{AD} = a\vect{AB} + b\vect{AC}$，则
由坐标得 $a - b = 1$，$b = 1$，$-a = 1$，前两式推出
$a = 2$，与 $a = -1$ 矛盾。四点不共面。

\emph{2.} $\vect{BC}(-2, 1, 1)$，$\vect{BD}(0, 1, 2)$。求解
$\vec n \cdot \vect{BC} = \vec n \cdot \vect{BD} = 0$，由
$b + 2c = 0$ 得 $b = -2c$；再由 $-2a - 2c + c = 0$ 得 $c = -2a$。取
$a = 1$，得 $\vec n(1, 4, -2)$。过点 $B(2,1,0)$ 的平面为：
\[
x + 4y - 2z - 6 = 0 .
\]

\emph{3.} 点 $A(1,1,1)$ 到平面的距离：
$\dfrac{\abs{1 + 4 - 2 - 6}}{\sqrt{21}} = \dfrac{3}{\sqrt{21}}$.
三角形 $BCD$ 的面积：$\norm{\vect{BC}}^2 = 6$，$\norm{\vect{BD}}^2 = 5$，
$\vect{BC}\cdot\vect{BD} = 3$，于是
\[
\text{面积} = \frac12\sqrt{6 \times 5 - 9} = \frac{\sqrt{21}}{2}.
\]

\emph{4.}
\[
V = \frac13 \times \frac{\sqrt{21}}{2} \times \frac{3}{\sqrt{21}}
= \frac12 .
\]
\end{solution}

\begin{solution}{pb:g12:space:1}
\textbf{1.} 平面：$x + y + z = 1$。直线：$(t, t, t)$。
交点满足 $3t = 1$：
$\left(\frac13, \frac13, \frac13\right)$.

\textbf{2.} $\dfrac{\abs{0 + 0 + 0 - 1}}{\sqrt3} =
\dfrac{1}{\sqrt3} = \dfrac{\sqrt3}{3}$.

\textbf{3.} $1 - 1 + 0 = 0$：正交。

\textbf{4.} 法向量相同，均为 $(1,1,1)$，但 $\frac52 \neq 1$：
两平面严格平行。直线的方向向量：
$(1,1,0) \cdot (1,1,1) = 2 \neq 0$，故不与平面平行，
因此它与平面相交于一点。

\textbf{5.} 各中点为 $\left(1, \frac12, 0\right)$，
$\left(\frac12, 1, 0\right)$, $\left(0, 1, \frac12\right)$,
$\left(0, \frac12, 1\right)$, $\left(\frac12, 0, 1\right)$,
$\left(1, 0, \frac12\right)$：每个点的坐标之和为
$\frac32$.

\textbf{6.} $P$ 的法向量为 $(1,1,1)$，
与 $(AG)$ 的方向一致，即直线垂直于平面。中心
$\left(\frac12,\frac12,\frac12\right)$ 的坐标之和
为 $\frac32$，故中心在 $P$ 上。

\textbf{7.} 相邻中点之差为形如
$\left(-\frac12, \frac12, 0\right)$ 的向量，长度为
$\frac{\sqrt2}{2}$，每次相同 --- 共六条相等的边。

\textbf{8.} 在 $\left(\frac12, 1, 0\right)$ 处，两条边向量为
$\left(\frac12, -\frac12, 0\right)$ 与
$\left(-\frac12, 0, \frac12\right)$，数量积为 $-\frac14$，
模均为 $\frac{\sqrt2}{2}$，$\cos = -\frac12$，夹角为
$120^\circ$。各边相等，各角均为 $120^\circ$，故为正
六边形。

\textbf{9.} 周长为 $6 \times \frac{\sqrt2}{2} = 3\sqrt2
\approx 4.24$。面积为 $6 \times \frac{\sqrt3}{4} \times
\frac12 = \frac{3\sqrt3}{4} \approx 1.30$，比一个面
（面积 $1$）还大 --- 这截面比正方体的任一面都大；只有
对角线矩形（$\sqrt2 \approx 1.41$）能超过它。

\textbf{10.} $(t,t,t)$ 与 $P$ 相交于 $t = \frac12$ 处，该点
即中心 --- 为 $[AG]$ 的中点（$A$ 对应 $t = 0$，$G$ 对应
$t = 1$）。

\textbf{11.} 当 $c = \frac12$ 时，平面截三条棱，
交点在 $A$ 所在的角上，分别为 $\left(\frac12,0,0\right)$，
$\left(0,\frac12,0\right)$、$\left(0,0,\frac12\right)$，截得一个
边长为 $\frac{\sqrt2}{2}$ 的正三角形。随着 $c$ 增大，
三角形变大，三个角被截去，形成六边形
（恰在 $c = \frac32$ 时为正），随后图形
对称地缩回 $G$ 所在顶点处的三角形：
这便是晶体切割式的形变。

\textbf{12.} 截面的每条边都是截平面与正方体的
一个\emph{面}的交集，而平面与一个面（平面多边形）至多
交于一条线段，故截面至多有
$6$ 条边。七条不可能；三至六条均可出现
（第 8 题与第 11 题即其中两例）。

\textbf{13.} 底面 $ABD$ 是面积为 $\frac12$ 的直角三角形；
高 $AE = 1$，体积为
$\frac13 \times \frac12 \times 1 = \frac16$.

\textbf{14.} $B, D, E, G$ 中任意两点恰在两个坐标上相差
$1$，六条距离均为 $\sqrt2$ --- 为正四面体。
正方体被分为 $BDEG$ 及四个角上的
与 $ABDE$ 全等的四面体，其体积为
$1 - 4 \times \frac16 = \frac13$.

\textbf{15.} 平面 $(BDE)$：$x + y + z = 1$；中心
$\left(\frac12,\frac12,\frac12\right)$ 到平面的距离为
$\frac{\abs{\frac32 - 1}}{\sqrt3} = \frac{1}{2\sqrt3} =
\frac{\sqrt3}{6}$.

\textbf{16.} $\vect{AG} = (1,1,1)$,
$\vect{BH} = (-1,1,1)$: $\cos\theta = \frac{-1 + 1 + 1}{3} =
\frac13$，$\theta \approx 70.5^\circ$，其补角为
$\approx 109.5^\circ$。四个价层电子对彼此排斥，
并在碳原子周围尽可能彼此远离：
最优的排布是正四面体的顶点，它们的
中心到顶点的方向恰好成这一角度 ---
金刚石就是第 14 题的结晶。

\textbf{17.} $(1, t, t)$ 在 $P$ 上：$1 + 2t = \frac32$，
故 $t = \frac14$，点为 $\left(1, \frac14, \frac14\right)$。

\textbf{18.} $(AB)$ 上的点为 $(t, 0, 0)$；$(EG)$ 上的点为
$(s, s, 1)$。若相交，则第三个坐标须
满足 $0 = 1$，这不可能。若平行，则需 $(1,0,0)$ 与
$(1,1,0)$ 共线，这不成立。既不相交也不平行，故两直线
为异面直线 --- 如同位于不同楼层的两条走廊。

\textbf{19.} $B + t(1,1,1) = (1 + t, t, t)$ 在 $P$ 上：
$1 + 3t = \frac32$，得 $t = \frac16$，投影点为
$\left(\frac76, \frac16, \frac16\right)$，其坐标之和
为 $\frac32$，故在 $P$ 上，符合要求。

\textbf{20.} 直线的参数式求交点（第 1、10、17 题）；
法向量用于度量角与距离（第 2、8、
15、16 题）；平面方程用于刻画截面（三角形到
六边形的形变）。这次游览收获颇丰：正方体中的正六边形、
角上的正四面体、金刚石的角度，以及两条彼此
互不相顾的直线 ---
只有空间才能提供的几何。
\end{solution}
